\documentclass[aspectratio = 43]{beamer}
\usetheme{AnnArbor}
% set captions with numbers
\setbeamertemplate{caption}[numbered]


\usepackage{amsmath, nccmath}
\usepackage[french]{babel}
\usepackage[latin1]{inputenc}
\usepackage[T1]{fontenc}


\definecolor{halfgray}{gray}{0.7}

\usecolortheme{whale}
\setbeamercolor{frametitle}{parent=structure,bg=halfgray}

% Include other packages here, before hyperref.
\DeclareGraphicsExtensions{.pdf,.png,.jpg,.jpeg}
\graphicspath{{Figures/}} %Where the figures folder is located

%%%%%%%%%%%%%%%%%%%%%%%%%%%%%%%%%%%%%%%%%%%%%%%%%%
\usepackage{pgffor}
\usepackage{listings}
\usepackage[%font=bf,
skip=0pt]{caption}

\captionsetup[lstlisting]{font={tiny,tt}}
\setbeamerfont{caption}{size=\tiny}
%\setlength\abovecaptionskip{-1pt}
\usepackage{pythonhighlight}
\usepackage{xcolor}

\definecolor{codegreen}{rgb}{0,0.6,0}
\definecolor{codegray}{rgb}{0.5,0.5,0.5}
\definecolor{codepurple}{rgb}{0.58,0,0.82}
\definecolor{backcolour}{rgb}{0.95,0.95,0.92}

\lstdefinestyle{mystyle}{%
    backgroundcolor=\color{backcolour},
    commentstyle=\color{codegreen},
    keywordstyle=\color{magenta},
    numberstyle=\tiny\color{codegray},
    stringstyle=\color{codepurple},
    basicstyle=\ttfamily\footnotesize,
    breakatwhitespace=false,
    breaklines=true,
    captionpos=b,
    keepspaces=true,
    numbers=left,
    numbersep=1pt,
    showspaces=false,
    showstringspaces=false,
    showtabs=false,
    tabsize=2,
}

\lstset{style=mystyle}
%\usepackage{python}
%%%%%%%%%%%%%%%%%%%%%%%%%%%%%%%%%%%%%%%%%%%%%%%%%%

\title{Revue de Projet}
\subtitle{Classification et Comptage de contenants vides}
\author[]{\tiny Thomas CHECCHIN $-$ Dorian CHEVALIERAS $-$ Nicolas TO VAN
  TRANG $-$ Za�d GHALI}
\institute[]{\textbf{T�l�com Physique Strasbourg}\\
\textbf{SEW - Usocome}}
\date{\tiny \today}


\begin{document}

%%%%--1--%%%%
\begin{frame}
  \begin{minipage}[t][0.2\textheight][t]{\textwidth}
    \centering
    \includegraphics[scale=0.3,width=0.2\textwidth]{tps-logo.png}
    \label{fig:tps_logo}
    \hspace*{4cm}
    \includegraphics[scale=0.3,width=0.2\textwidth]{SEW_logo.jpg}
    \label{fig:sew_logo}
  \end{minipage}
  \titlepage%
\end{frame}
%%%% --2--%%%%
\begin{frame}
  \frametitle{Plan}
  \tableofcontents
\end{frame}
%
% %%%%--3--%%%%
\section{Remise en contexte}
\begin{frame}
  \frametitle{Remise en contexte}
  \begin{itemize}
  \item Pr�sentation entreprise
  \item Etat actuel
  \end{itemize}
\end{frame}
%
% %%%% --4--%%%%
\section{Pr�vision des t�ches � r�aliser}
\begin{frame}
  \frametitle{Pr�vision des t�ches � r�aliser}
  \begin{itemize}
  \item Pr�sentation du cahier des charges + type de cam�ra envisag�
    donc budgets ?
  \item Diagramme de Gantt sur l'enti�ret� du projet
  \item Diagramme de Gantt pr�visionnel avant le R1 et diagramme r�el
  \end{itemize}
\end{frame}
%
% %%%% --5--%%%%
\section{Etat de l'art}
\begin{frame}
  \frametitle{Remise en contexte}
  \begin{itemize}
  \item rfid (trop cher)
  \item vision (contours des bo�tes)
  \item Plus important : YOLO mais d'autres existent mais moins
    utilis�s comme \textit{SDD ?}
  \end{itemize}
\end{frame}
%
% %%%% --6--%%%%
\section{Pistes de solutions}
\begin{frame}
  \frametitle{Pistes de solutions}
  \begin{itemize}
  \item Pour les zones de production : comptage apr�s identification
    des bo�tes vide ou non en ayant ou non d�finit quel type de bo�te
  \item Pour la zone de stockage : utilisation des zones strat�giques
    de passages pour une meilleure classification et s'affranchir des
    risques d'identification des bo�tes cach�es
  \end{itemize}
\end{frame}
%
% %%%% --7--%%%%
\section{Conclusion}
\begin{frame}
  \frametitle{Conclusion}
  Poursuite du PI ?
\end{frame}
\end{document}
